% 詹姆斯·格雷果里（综述）
% license CCBYSA3
% type Wiki

本文根据 CC-BY-SA 协议转载翻译自维基百科\href{https://en.wikipedia.org/wiki/James_Gregory_(mathematician)}{相关文章}

詹姆斯·格雷戈里（James Gregory FRS，1638年11月－1675年10月）是一位苏格兰数学家和天文学家。他的姓氏有时拼作Gregorie，这是苏格兰的原始拼法。他提出了反射望远镜的早期实用设计——格里高利望远镜，并在三角学方面取得进展，发现了若干三角函数的无穷级数展开。

在他的著作《几何学的普遍部分》（Geometriae Pars Universalis，1668年）[1] 中，格雷戈里首次发表并证明了微积分基本定理（从几何角度陈述，并且仅限于后来版本定理所考虑曲线的一类特殊情况）。艾萨克·巴罗对此予以承认。[2][3][4][5][6][7]
